% =====================================================================
%  Sezione 2.6 — Grandezze integrali: accumulare invece di derivare
%  Segue la sezione sulle tre rappresentazioni geometriche (A1, A2, A3)
%  e precede quella sull'altezza dal pavimento.
% =====================================================================

\section{Grandezze integrali: accumulare invece di derivare}
\label{sec:integrali}

Le tre rappresentazioni descritte finora contengono, accanto alle grandezze
statiche, alcune grandezze dinamiche: velocità del bacino, velocità della testa,
variazione dell'angolo del tronco. Sono le grandezze che dovrebbero distinguere
una caduta da un cambio di postura controllato, perché la differenza fra i due
eventi è per sua natura dinamica.

Questa sezione mostra che, calcolate come derivate istantanee, quelle grandezze
non svolgono il ruolo che dovrebbero, e propone una formulazione alternativa.

\subsection{Il modello non usava la dinamica}

L'osservazione da cui parte l'analisi è sperimentale. Verificando il
comportamento del classificatore sulle registrazioni domestiche, è emerso che
esso assegnava punteggio elevato anche a corpi \emph{completamente immobili}, e
che la correlazione fra il punteggio prodotto e il movimento effettivamente
presente nella finestra era pari a $+0{,}139$, cioè praticamente nulla.

Il modello, in altre parole, aveva imparato una scorciatoia posturale: un tronco
orizzontale implicava una caduta, indipendentemente da come il corpo fosse
arrivato in quella configurazione. È una regola perfetta su un dataset in cui
nessuno si sdraia mai su un divano, e sbagliata in una casa vera.

La prima ipotesi era che il difetto fosse del classificatore. La verifica ha
mostrato che il difetto era \emph{a monte}, nelle grandezze che gli venivano
fornite.

\subsection{Perché derivare amplifica il rumore}

I \emph{landmark} non sono misure dirette del corpo ma stime prodotte dal
\emph{pose estimator}, e come tali sono affette da rumore. Una velocità
calcolata come differenza fra fotogrammi consecutivi,

\[
v(t)=\frac{x(t)-x(t-1)}{\Delta t},
\]

divide la differenza fra \emph{due stime rumorose} per un intervallo di tempo
piccolo. Se ciascuna stima ha un errore di deviazione standard $\sigma$, la
differenza ne ha $\sigma\sqrt{2}$, e la divisione per $\Delta t = 1/\text{fps}$
la moltiplica per il frame rate. A 18 fotogrammi al secondo, l'errore sulla
velocità è quindi circa venticinque volte quello sulla posizione.

La derivazione, in sostanza, è un'operazione che amplifica le alte frequenze —
ed è esattamente lì che vive il rumore dello stimatore di posa.

\subsection{Una misura del fenomeno: l'autocorrelazione a un passo}

Per quantificare quanto di ciascuna grandezza sia segnale e quanto rumore è
stata misurata l'\textbf{autocorrelazione a un passo}, cioè la correlazione fra
il valore della grandezza al fotogramma $t$ e il suo valore al fotogramma
$t-1$:

\[
\rho_1 = \frac{\operatorname{Cov}\big(x_t,\,x_{t-1}\big)}
              {\sigma_{x_t}\,\sigma_{x_{t-1}}}.
\]

L'idea è che un segnale fisico vero cambia con continuità e quindi è molto
correlato con sé stesso da un fotogramma al successivo; un rumore bianco non lo
è affatto. I valori misurati sono riportati in
Tabella~\ref{tab:autocorrelazione}.

\begin{table}[ht]
\centering
\caption{Autocorrelazione a un passo delle grandezze dinamiche, misurata sui
fotogrammi con posa valida di UP-Fall.}
\label{tab:autocorrelazione}
\begin{tabular}{lc}
\hline
\textbf{Grandezza} & $\boldsymbol{\rho_1}$ \\
\hline
\multicolumn{2}{l}{\emph{derivate istantanee}} \\
\quad Variazione dell'angolo del tronco & 0,176 \\
\quad Energia del corpo & 0,228 \\
\quad Velocità del bacino & 0,322 \\
\quad Velocità della testa & 0,438 \\
\hline
\multicolumn{2}{l}{\emph{grandezze integrali}} \\
\quad Escursione del tronco & 0,916 \\
\quad Percorso del bacino & 0,939 \\
\quad Discesa del bacino & 0,911 \\
\hline
\end{tabular}
\end{table}

Il contrasto è netto: le derivate stanno fra $0{,}18$ e $0{,}44$, le grandezze
integrate fra $0{,}91$ e $0{,}94$.

È però necessaria una precisazione, perché l'autocorrelazione non è una prova
sufficiente. Un valore basso non dimostra da solo che una grandezza sia rumore,
così come un valore alto non dimostra che sia informativa per il
riconoscimento; e una grandezza ottenuta integrando su un intervallo tende
\emph{per costruzione} a essere più regolare nel tempo. L'autocorrelazione va
quindi usata come \textbf{spiegazione del fenomeno}, non come sua
dimostrazione. La dimostrazione è il comportamento discriminativo, riportato
nel Capitolo~\ref{cap:sperimentale}.

\subsection{La formulazione integrale}

L'alternativa alla derivazione è l'\textbf{accumulazione}: invece di misurare
quanto il tronco ruota \emph{adesso}, si misura di quanto ha ruotato \emph{in
tutto} nell'ultimo mezzo secondo. È una descrizione più fedele al fenomeno,
perché una caduta \emph{è} un'escursione, e produce grandezze che variano
lentamente invece di oscillare a ogni fotogramma.

È necessario però dichiarare subito che cosa questa formulazione \textbf{non}
fa, perché l'intuizione naturale è sbagliata e conduce a una conclusione
errata. Si sarebbe tentati di affermare che accumulando il rumore si cancella,
poiché un errore a media nulla sommato su molti fotogrammi tende a zero.
L'argomento vale per una somma di incrementi \emph{con segno}, ma le tre
grandezze qui introdotte non lo sono:

\begin{itemize}
    \item il \emph{percorso} è una somma di valori assoluti, e i valori assoluti
          non si cancellano mai;
    \item l'\emph{escursione} è una differenza fra massimo e minimo, e il rumore
          allarga l'intervallo;
    \item la \emph{discesa} è uno scarto dal massimo recente, e il rumore
          abbassa il minimo.
\end{itemize}

Tutte e tre sono quindi funzioni \textbf{monotone crescenti} del rumore. Ciò che
le rende preferibili alle derivate non è la robustezza, ma la
\textbf{regolarità}: variano lentamente per costruzione, e non presentano le
discontinuità che rendono le derivate quasi indistinguibili dal rumore di
fondo. La Sezione~\ref{sec:rumore-int} misura fin dove questa distinzione
regge.

Sono state quindi introdotte tre grandezze integrali, calcolate su una finestra
mobile di $K$ fotogrammi e mai divise per $\Delta t$. Detta $\theta_t$ la
sequenza dell'angolo del tronco e $(x_t, y_t)$ quella della posizione del punto
seguito, normalizzata da una scala $s$ che rende le grandezze adimensionali:

\begin{align}
\text{escursione}_t &= \max_{\tau \in W} \theta_\tau
                     - \min_{\tau \in W} \theta_\tau
                     \label{eq:escursione}\\[4pt]
\text{percorso}_t   &= \frac{1}{s}\sum_{\tau \in W}
                       \sqrt{(x_\tau - x_{\tau-1})^2 + (y_\tau - y_{\tau-1})^2}
                     \label{eq:percorso}\\[4pt]
\text{discesa}_t    &= \frac{1}{s}\Big( y_t - \min_{\tau \in W} y_\tau \Big)
                     \label{eq:discesa}
\end{align}

dove $W = \{t-K+1,\dots,t\}$ è la finestra di accumulazione. Le tre grandezze
rispondono a tre domande diverse: di quanto il tronco ha ruotato in tutto,
quanta strada ha percorso il corpo, e di quanto è sceso rispetto al punto più
alto raggiunto di recente.

La normalizzazione $s$ è la lunghezza del tronco per le rappresentazioni
riferite al corpo e l'asse maggiore dell'ellisse per quella basata su PCA:
senza di essa il percorso sarebbe espresso in pixel e dipenderebbe dalla
distanza dalla telecamera.

\subsection{Il vincolo di causalità}

Una precisazione implementativa che ha conseguenze pratiche importanti.

Le equazioni~\eqref{eq:escursione}--\eqref{eq:discesa} si possono calcolare in
due modi. Il primo è applicare una funzione di finestra mobile a un file di
feature già completo, quindi \emph{a registrazione conclusa}. È il modo più
semplice, e produce esattamente gli stessi numeri per un'analisi fuori linea.

Il secondo è mantenere un \textbf{buffer circolare} degli ultimi $K$
fotogrammi e aggiornarlo a ogni nuovo campione. È il modo più laborioso, ed è
l'unico compatibile con il funzionamento su dispositivo: la board riceve un
fotogramma e deve decidere, il video intero non lo possiede.

La distinzione non è formale. In una versione precedente di questo lavoro le
grandezze integrali erano calcolate nel primo modo, ed è la ragione per cui non
sono mai arrivate sul dispositivo: i risultati che le riguardavano erano tutti
calcoli fuori linea, mentre il sistema installato utilizzava la sola
rappresentazione senza integrali. Nell'implementazione attuale il motore
geometrico mantiene il buffer, e le stesse grandezze si calcolano identiche
fuori linea e in tempo reale.

Il costo è un ritardo di riscaldamento: nei primi $K$ fotogrammi dopo l'avvio
il buffer non è pieno e i valori non sono affidabili. Con $K = 31$ a 18
fotogrammi al secondo, il periodo di riscaldamento è di circa 1,7 secondi, ed è
dichiarato esplicitamente dal sistema.

\subsection{La scelta di $K$}

La lunghezza della finestra di accumulazione regola un compromesso: un valore
piccolo mantiene la reattività ma cancella poco rumore, un valore grande
stabilizza il segnale ma diluisce l'evento in fotogrammi normali.

Il valore adottato è $K = 31$ fotogrammi, pari a circa 1,7 secondi a 18
fotogrammi al secondo. La scelta non è arbitraria: nel
Capitolo~\ref{cap:sperimentale} è riportata la scansione su cinque valori, dalla
quale risulta che il parametro sposta il tasso di falsi allarmi in ambiente
reale di un fattore undici. È il tipo di parametro che, non dichiarato, rende
un risultato non riproducibile.

\subsection{Fin dove regge: rumore indipendente e ambiguità di posa}
\label{sec:rumore-int}

Poiché le tre grandezze crescono con il rumore, è necessario misurare quanto
questa proprietà le penalizzi in pratica. La verifica è stata condotta
iniettando rumore gaussiano indipendente sulle coordinate dei \emph{landmark}
di UP-Fall, con deviazione espressa in frazione della lunghezza del torso, e
rimisurando la prestazione.

\begin{table}[ht]
\centering
\caption{F2 al crescere del rumore iniettato sui \emph{landmark}. Fra parentesi
la variazione rispetto all'assenza di rumore.}
\label{tab:rumore}
\begin{tabular}{lccc}
\hline
$\boldsymbol{\sigma}$ & \textbf{senza integrali} & \textbf{con integrali} & \textbf{sottoinsieme scelto} \\
\hline
0     & 0,874 & 0,858 & 0,857 \\
0,005 & 0,863 \,(--0,011) & 0,864 \,(+0,006) & 0,844 \,(--0,012) \\
0,010 & 0,855 \,(--0,018) & 0,866 \,(+0,008) & 0,863 \,(+0,007) \\
0,020 & 0,864 \,(--0,010) & 0,846 \,(--0,012) & 0,836 \,(--0,021) \\
0,040 & 0,863 \,(--0,011) & 0,856 \,(--0,002) & 0,825 \,(--0,031) \\
\hline
\end{tabular}
\end{table}

Il rumore indipendente \textbf{non degrada apprezzabilmente} nessuno dei tre
insiemi: quadruplicando l'ampiezza rispetto a quella stimata sulle
registrazioni domestiche, la perdita massima è $0{,}031$, inferiore alla
dispersione fra soggetti ($\pm 0{,}08$). Le grandezze integrali, in
particolare, non si degradano più delle derivate.

La spiegazione è che un errore indipendente su ciascun punto e a ciascun
fotogramma si media comunque, anche attraverso una funzione monotona, perché
punti diversi sbagliano in direzioni diverse.

\subsubsection*{L'ambiguità di posa}

Il rumore indipendente non è però l'unica modalità di errore dello stimatore,
e non è quella che si osserva sul campo. In una ripresa domestica in cui il
soggetto è vicino alla telecamera e con gli arti scorciati verso l'obiettivo,
lo stimatore non produce oscillazioni casuali: \textbf{alterna fra due
interpretazioni plausibili della stessa configurazione} --- quale arto sia
davanti, dove cada il bacino --- e ciascuna interpretazione resta stabile per
diversi fotogrammi consecutivi.

L'effetto sulle grandezze accumulate è qualitativamente diverso, ed è stato
misurato:

\begin{table}[ht]
\centering
\caption{Escursione del tronco su finestra mobile, in riprese domestiche.}
\label{tab:ambiguita}
\begin{tabular}{lcc}
\hline
 & \textbf{mediana} & \textbf{massimo} \\
\hline
Persona \textbf{ferma} su un divano & 9,0$^\circ$ & \textbf{93,5}$^\circ$ \\
\textbf{Caduta} reale & 9,0$^\circ$ & 60,7$^\circ$ \\
\hline
\end{tabular}
\end{table}

Una persona immobile produce un'escursione superiore a quella di una caduta.
Non si tratta di movimento: è la stima che salta fra due configurazioni, e per
un accumulatore un salto coerente ripetuto è indistinguibile da un movimento
reale.

Questa è dunque la vulnerabilità effettiva della formulazione integrale, ed è
diversa da quella che si sarebbe supposta: non il rumore, ma l'ambiguità. La
difesa non può essere un filtro, perché il segnale spurio è regolare quanto
quello vero; deve essere un criterio di \textbf{affidabilità della stima}. Lo
stimatore fornisce una confidenza per ciascun punto --- la stessa grandezza che
ha rivelato l'occlusione del volto discussa nel
Capitolo~\ref{cap:sperimentale} --- e quando la confidenza dei punti utilizzati
cade durante la finestra di accumulazione l'accumulo non è affidabile: il
sistema deve dichiararlo invece di decidere, come già fa durante il periodo di
riscaldamento.

\subsection{Cosa sostituiscono}

L'introduzione delle grandezze integrali si accompagna alla rimozione di tre
grandezze presenti nella formulazione iniziale, denominate \emph{energie} e
calcolate come somme di quadrati dei contributi dinamici di ciascuna
rappresentazione.

Il motivo della rimozione è dimensionale: quelle somme sommavano quadrati di
gradi, di rapporti adimensionali, di pixel e di frazioni di immagine, cioè
quantità non omogenee. Il risultato non ha un'unità di misura interpretabile, e
sul dataset presentava rapporti fra valore massimo e mediano dell'ordine di
$7\cdot10^{5}$. Erano inoltre all'origine della disparità di scala fra le
grandezze in ingresso, che senza normalizzazione rendeva alcune colonne
invisibili al modello.

Le grandezze integrali svolgono lo stesso ruolo — riassumere il contenuto
dinamico di una finestra in un numero — ma sono dimensionalmente coerenti,
normalizzate rispetto a una scala corporea e interpretabili singolarmente.
